\section{Hospitais que não são comparáveis entre si: a reestimação
segmentada}
\label{sec:segmentacao}

Todos os resultados até aqui compararam hospitais de uma rede com
hospitais de outra, descontando complexidade e porte por regressão. Esta
seção faz uma pergunta anterior a essa: \textbf{há hospitais no painel que
não deveriam estar sendo comparados de jeito nenhum?}

\subsection{O caso que motivou a pergunta}

O gatilho foi concreto e está na Seção~\ref{sec:categorias}. A ocupação de
UTI só é definida onde há leito de UTI, e medi-la sobre a rede inteira
somaria hospitais \textbf{que não têm UTI} a hospitais de UTI vazia, que são
coisas diferentes.

Ali o problema é de definição, mas ele expõe um problema mais geral. Um
hospital de pequeno porte do interior, sem UTI e sem alta complexidade, e um
hospital-escola de mil leitos com UTI e transplante não fazem a mesma coisa.
Medi-los pelos mesmos seis indicadores e descontar a diferença com uma
covariável de complexidade supõe que a diferença entre eles é de
\emph{grau}. Boa parte dela é de \emph{gênero}: há indicadores que
simplesmente não se aplicam a um dos dois.

A resposta natural seria cruzar duas dimensões, ter ou não UTI e fazer ou
não alta complexidade, e obter quatro grupos. \textbf{Esse cruzamento
degenera.} A célula ``sem UTI, mas faz alta complexidade'' é
\textbf{vazia}: no painel de 289 hospitais não existe um único caso. Alta
complexidade sem terapia intensiva não é uma combinação rara: é uma
combinação que a organização assistencial não produz.

O que sobra, portanto, não são quatro grupos: é uma \textbf{escada de três
degraus}, ordenada por capacidade instalada.

\subsection{Os três estratos, e por que ficam fixos}

\begin{itemize}
\item \textbf{E1, sem UTI SUS.} Mediana de leitos de UTI SUS igual a zero
  nos onze anos.
\item \textbf{E2, com UTI, baixa complexidade.} Mediana de UTI positiva e
  faixa Barcelona abaixo de 4.
\item \textbf{E3, com UTI, alta complexidade.} Mediana de UTI positiva e
  faixa Barcelona 4 ou mais.
\end{itemize}

A classificação é \textbf{fixa para todo o período}, e essa escolha precisa
de justificativa porque um hospital pode, em princípio, mudar de degrau ao
longo de onze anos. Duas verificações sustentam a fixação.

A primeira é direta: \textbf{a faixa Barcelona não varia dentro de nenhum
CNES do painel}. Ela é atributo estrutural, e não há o que fixar: já é
fixa.

A segunda é menos óbvia e mais interessante. O número de leitos de UTI
\emph{oscila} entre anos, e à primeira vista isso pediria classificação
anual. Mas a oscilação não é obra física: é \textbf{suspensão e limpeza de
cadastro}, tipicamente inativação temporária por falta de intensivista. A
evidência é que existem 38 observações, em 26 hospitais, com \textbf{zero
leito de UTI cadastrado e diárias de UTI registradas no mesmo ano}: a ala
operando enquanto o registro está zerado. O controle equivalente na
internação, leito zerado com diária registrada, dá \textbf{zero casos}, como
se espera de um leito comum, que não é habilitação que se suspenda. Sendo o
cadastro que pisca e não a capacidade que muda, \textbf{a mediana dos onze
anos é o estimador certo} da capacidade subjacente: ela ignora os anos em
que o registro caiu sem que a ala fechasse.

\subsection{O funil}

\begin{table}[H]
\centering
\caption{Distribuição dos 289 hospitais pelos três estratos estruturais,
por categoria de gestão. Cada célula traz hospitais distintos (CNES), e
\textbf{os cinco conversores são contados nas duas categorias que ocuparam},
pela mesma convenção declarada na Seção~\ref{sec:categorias}. Por isso a
soma das células de uma coluna excede o total dela: em E2 por um hospital
(Presidente Prudente) e em E3 por quatro (Sorocaba, Pérola Byington, Taipas
e Vila Penteado). Os totais contam hospitais distintos, uma vez cada.}
\label{tab:funilestratos}
\small
\begin{tabular}{lrrrr}
\toprule
Categoria & E1 (sem UTI) & E2 (UTI, baixa) & E3 (UTI, alta) & Total \\
\midrule
Direta            &  1 &  6 &  20 &  27 \\
OSS               &  2 &  7 &  28 &  37 \\
Público Municipal &  9 & 16 &  31 &  56 \\
Filantrópico      & 54 & 67 &  50 & 171 \\
Privado           &  1 &  2 &   0 &   3 \\
\midrule
\textbf{Total}    & \textbf{67} & \textbf{97} & \textbf{125} & \textbf{289} \\
\bottomrule
\end{tabular}
\end{table}

A tabela já entrega o principal achado estrutural da segmentação, e ele é
sobre \textbf{quem pode ser comparado com quem}. Em E1 e E2, a administração
Direta tem 1 e 6 hospitais, e as OSS têm 2 e 7. \textbf{Nesses dois estratos
não há base para comparação entre redes}, e é por isso que os resultados de
E1 e E2 aparecem tarjados em todas as tabelas desta seção.

É só em \textbf{E3} que existe comparação viável: \textbf{28 hospitais OSS
contra 20 da administração Direta}. E é também em E3 que estão quatro dos
cinco hospitais que trocaram de rede no período.

\subsection{A segmentação captura gradiente estrutural real}

Antes de usar os estratos, vale verificar se eles significam alguma coisa. A
resposta é sim: \textbf{em cinco dos seis indicadores, as medianas ordenam
monotonicamente}, E1, depois E2, depois E3, sem inversões. Em quatro
deles a ordenação atravessa os três estratos; no quinto, a ocupação de
UTI, ela compara apenas E2 e E3, pelo motivo registrado logo abaixo da
tabela.

\begin{table}[H]
\centering
\caption{Mediana de cada indicador por estrato, todas as categorias
reunidas. A ordenação é sempre pela mediana.}
\label{tab:gradienteestratos}
\small
\begin{tabular}{lrrrl}
\toprule
Indicador & E1 & E2 & E3 & Ordena? \\
\midrule
Ocupação de internação (\%)   & 41,0   & 66,1   & 90,2   & sim \\
Faturamento real por saída (R\$) & 658    & 1.294  & 1.821  & sim \\
TMP (dias)                    & 3,19   & 3,84   & 5,25   & sim \\
Fração de alta complexidade   & 0,000  & 0,001  & 0,054  & sim \\
Ocupação de UTI (\%)          & ---    & 76,1   & 77,3   & sim \\
Mortalidade geral             & 0,0438 & 0,0585 & 0,0587 & \textbf{empata} \\
\bottomrule
\end{tabular}
\end{table}

A exceção é a \textbf{mortalidade}, e ela é informativa. Entre E2 e E3 as
medianas são 0,0585 e 0,0587: \textbf{empate}. A mortalidade separa bem
``sem UTI'' de ``com UTI'', mas não distingue baixa de alta complexidade.
Isto é coerente com o que se espera: a presença de terapia intensiva muda o
perfil de gravidade dos pacientes atendidos; a faixa de complexidade, dentro
de quem já tem UTI, muda mais o \emph{tipo} de procedimento do que o risco
de morte.
\newpage
\textbf{Sobre a célula vazia da ocupação de UTI em E1, que não é bem uma
célula vazia.} E1 é definido pela \emph{mediana} dos onze anos de leitos de
UTI SUS igual a zero, e não por ausência de UTI em todos os anos. Dos 67
hospitais do estrato, \textbf{18 registram UTI SUS em algum ano}, num total
de 71 hospitais-ano, e para eles a razão existe e vale 66,2\% na mediana.
São o \emph{piscar cadastral} descrito acima: a ala opera e o registro cai.
A célula fica fora da tabela porque 71 observações em 18 hospitais não
descrevem o estrato, \textbf{não porque a razão seja indefinida ali}.

\comoler{nos boxplots, procure se as três caixas de cada figura sobem da
esquerda para a direita sem se cruzar; é isso que significa ``ordena''. Na
mortalidade, as caixas de E2 e E3 ficam praticamente na mesma altura.}

\comofeito{boxplots por indicador: o eixo horizontal traz os três estratos
estruturais (E1, E2 e E3), subdivididos pelas cinco categorias de gestão, e
o vertical traz o valor do indicador; a linha dentro da caixa é a mediana e
a caixa vai do primeiro ao terceiro quartil. O asterisco marca a célula com
menos de dez hospitais, que não sustenta leitura. A unidade é o
hospital-ano.}

\begin{figure}[H]
\centering
\includegraphics[width=\textwidth]{figSEG_01_box_ocupacao_internacao.png}
\caption{Ocupação de internação por estrato estrutural e categoria de
gestão. É a figura em que a distância entre as redes mais aparece: dentro
do estrato E3, o único com hospitais dos dois lados em número suficiente, a
caixa das OSS está inteira acima da caixa da administração Direta.}
\label{fig:estratos}
\end{figure}

\begin{figure}[H]
\centering
\includegraphics[width=\textwidth]{figSEG_01_box_tmp.png}
\caption{Tempo médio de permanência por estrato estrutural e categoria de
gestão.}
\label{fig:estratostmp}
\end{figure}
\newpage
\begin{figure}[H]
\centering
\includegraphics[width=\textwidth]{figSEG_01_box_mort_all.png}
\caption{Mortalidade geral por estrato estrutural e categoria de gestão.}
\label{fig:estratosmort}
\end{figure}

\begin{figure}[H]
\centering
\includegraphics[width=\textwidth]{figSEG_01_box_custo_real.png}
\caption{Faturamento real por saída (R\$ de dez/2025) por estrato
estrutural e categoria de gestão, em escala logarítmica.}
\label{fig:estratoscusto}
\end{figure}
\newpage
\begin{figure}[H]
\centering
\includegraphics[width=\textwidth]{figSEG_01_box_ocupacao_uti_valida.png}
\caption{Ocupação de UTI por estrato estrutural e categoria de gestão,
calculada apenas onde há leito de UTI SUS cadastrado.}
\label{fig:estratosuti}
\end{figure}

\begin{figure}[H]
\centering
\includegraphics[width=\textwidth]{figSEG_01_box_pct_alta_complex.png}
\caption{Fração de alta complexidade por estrato estrutural e categoria de
gestão.}
\label{fig:estratosac}
\end{figure}
\newpage
\comoler{nas séries, o que importa é se as três linhas correm separadas a
década inteira ou se encostam. Onde correm separadas, o estrato é uma
divisão estável da rede e não um retrato de um ano.}

\comofeito{o eixo horizontal traz o ano (2015 a 2025) e o vertical, a
mediana do indicador naquele ano, uma linha por estrato. A unidade é o
hospital-ano, e todas as categorias de gestão entram juntas em cada
estrato. A faixa cinza marca 2020 e 2021. A curva de E1 na ocupação de UTI
é rala e sai tracejada, porque ali só existem os 71 hospitais-ano do
piscar cadastral.}

\begin{figure}[H]
\centering
\includegraphics[width=\textwidth]{figSEG_02_serie_ocupacao_internacao.png}
\caption{Ocupação de internação: mediana anual por estrato estrutural.}
\label{fig:estratosserieocup}
\end{figure}

\begin{figure}[H]
\centering
\includegraphics[width=\textwidth]{figSEG_02_serie_pct_alta_complex.png}
\caption{Fração de alta complexidade: mediana anual por estrato
estrutural. As linhas de E1 e E2 correm coladas ao eixo a década inteira:
alta complexidade, na rede paulista, é assunto de E3.}
\label{fig:estratosserieac}
\end{figure}
\newpage
\begin{figure}[H]
\centering
\includegraphics[width=\textwidth]{figSEG_02_serie_tmp.png}
\caption{Tempo médio de permanência: mediana anual por estrato estrutural.}
\label{fig:estratosserietmp}
\end{figure}

\begin{figure}[H]
\centering
\includegraphics[width=\textwidth]{figSEG_02_serie_mort_all.png}
\caption{Mortalidade geral: mediana anual por estrato estrutural.}
\label{fig:estratosseriemort}
\end{figure}
\newpage
\begin{figure}[H]
\centering
\includegraphics[width=\textwidth]{figSEG_02_serie_custo_real.png}
\caption{Faturamento real por saída (R\$ de dez/2025): mediana anual por
estrato estrutural, em escala logarítmica.}
\label{fig:estratosseriecusto}
\end{figure}

\begin{figure}[H]
\centering
\includegraphics[width=\textwidth]{figSEG_02_serie_ocupacao_uti_valida.png}
\caption{Ocupação de UTI: mediana anual por estrato estrutural, calculada
apenas onde há leito de UTI SUS cadastrado.}
\label{fig:estratosserieuti}
\end{figure}

\subsection{O que os estratos separam, e o que eles não separam}
\label{sub:ampliada}

A tabela anterior responde a uma pergunta binária: a mediana ordena ou não.
Ela não diz \emph{quanto} os estratos separam, nem em quais dimensões
separam muito e em quais quase nada. Esta subseção responde isso cruzando os
três estratos com \textbf{as 27 variáveis do painel que admitem leitura
descritiva}, e não apenas com os seis indicadores do estudo.

\subsubsection*{Duas decisões que precisam vir antes dos números}

\textbf{Primeira: a unidade é o hospital, não o hospital-ano.} O estrato é
fixo por CNES, de modo que as onze observações de um hospital são o mesmo
hospital medido onze vezes. Para \emph{descrever} os estratos o hospital-ano
serve, e é o que a tabela anterior usa; para \emph{comparar} os estratos ele
não serve, porque trata como independentes onze medidas do mesmo objeto e
infla qualquer teste por um fator próximo de onze. Os testes desta subseção
usam uma linha por hospital (a mediana dos onze anos), com $n$ igual a 289.

\textbf{Segunda: sete das 27 variáveis são definicionais e não contam.} O
corte em E1, E2 e E3 é construído a partir da mediana de leitos de UTI SUS e
da faixa Barcelona. Encontrar que os estratos separam bem a faixa Barcelona
($\eta^2$ de 0,71) ou a proporção de leitos de UTI ($\eta^2$ de 0,54) não é
achado: é a definição do corte reaparecendo. Essas sete estão marcadas na
figura e na tabela-fonte, e ficam fora de toda leitura substantiva.

\begin{figure}[H]
\centering
\includegraphics[width=.94\textwidth]{figSEG_11_eta2_ranking.png}
\caption{Parcela da variância \emph{entre hospitais} que o estrato explica,
em cada uma das 27 variáveis. As barras hachuradas em cinza são as
definicionais: o corte é construído com elas.}
\label{fig:eta2estratos}
\end{figure}

\comoler{quanto mais comprida a barra, mais o perfil de complexidade do
hospital determina aquela variável. Leia de cima para baixo pulando as
barras hachuradas, que são a definição do corte. O que interessa é o
contraste entre o topo e o fundo da lista.}

\comofeito{gráfico de barras horizontais. Cada barra é uma variável do
painel, e o comprimento é o $\eta^2$: a fatia da diferença entre hospitais
que o degrau da escada explica sozinho, calculada sobre uma linha por
hospital (mediana dos onze anos, $n$ igual a 289). A cor indica se essa
diferença entre os três grupos é grande demais para ser acaso, pelo mesmo
teste usado na tabela adiante. As barras hachuradas são as variáveis com
que o próprio corte é construído.}

\subsubsection*{O achado central: o gradiente é estrutural, quase não é clínico}

A tabela a seguir tem duas colunas que precisam de tradução, e elas são
simples de entender uma vez ditas em português.

A coluna \textbf{$\eta^2$} responde: \emph{de tudo o que faz um hospital ser
diferente do outro neste indicador, quanto o degrau da escada já explica
sozinho?} Ela vai de 0 a 1. Um $\eta^2$ de 0,43 na ocupação quer dizer que
saber apenas em qual degrau o hospital está já explica 43\% da diferença de
ocupação entre dois hospitais quaisquer da rede --- e que os outros 57\%
dependem de tudo o mais: gestão, região, perfil de paciente, sorte. Quanto
mais alto, mais o degrau \emph{define} o hospital naquele indicador.

A coluna \textbf{$p_{23}$} responde outra coisa: \emph{a diferença entre o
segundo e o terceiro degrau é real ou pode ser acaso?} Ela é o resultado de
um teste que compara os dois grupos e devolve um número entre 0 e 1. Perto
de zero significa ``esta diferença é grande demais para ser coincidência'';
acima de 0,05, significa ``os dois grupos são indistinguíveis com os dados
que temos''. O teste usado compara a posição dos hospitais na fila em vez
das médias, o que é o certo aqui porque vários destes indicadores têm caudas
longas que distorceriam a média. E, como a tabela faz três comparações de
uma vez (E1 contra E2, E1 contra E3, E2 contra E3), o número foi ajustado
para levar isso em conta: quem faz várias comparações acaba encontrando
alguma ``diferença'' por puro acaso, e o ajuste desconta esse efeito. A
última coluna, \textbf{Ordena?}, resume tudo num sim ou não: os três degraus
sobem na ordem esperada e sem empate?

\begin{table}[H]
\centering
\caption{As duas pontas do ordenamento, entre as 20 variáveis não
definicionais. Os valores são medianas por hospital. $\eta^2$ é a fatia da
diferença entre hospitais que o degrau da escada explica sozinho, e
$p_{23}$ mede se a diferença entre o segundo e o terceiro degrau é real ou
pode ser acaso, com o ajuste para as três comparações feitas.}
\label{tab:eta2estratos}
\small
\begin{tabular}{lrrrrrl}
\toprule
Variável & E1 & E2 & E3 & $\eta^2$ & $p_{23}$ & Ordena? \\
\midrule
\multicolumn{7}{l}{\emph{o que o estrato explica}}\\
Ocupação de internação (\%)      & 39,2 & 66,0 & 90,6 & 0,43 & $<$0,001 & sim \\
Saídas hospitalares              & 2.211 & 4.266 & 10.460 & 0,43 & $<$0,001 & sim \\
Diárias de internação            & 7.179 & 16.450 & 53.200 & 0,41 & $<$0,001 & sim \\
Leitos SUS                       & 53 & 71 & 175 & 0,33 & $<$0,001 & sim \\
Faturamento real total (R\$ mil) & 1.341 & 5.295 & 18.900 & 0,23 & $<$0,001 & sim \\
TMP (dias)                       & 3,15 & 3,80 & 5,27 & 0,16 & $<$0,001 & sim \\
\midrule
\multicolumn{7}{l}{\emph{o que ele quase não explica}}\\
Faturamento real por saída (R\$) & 647 & 1.297 & 1.821 & 0,10 & $<$0,001 & sim \\
Leitos SUS como \% do total      & 80,5 & 63,7 & 81,2 & 0,08 & 0,001 & \textbf{não} \\
Fração de alta complexidade      & 0,000 & 0,001 & 0,054 & 0,07 & $<$0,001 & sim \\
\textbf{Mortalidade geral}       & 0,043 & 0,059 & 0,058 & \textbf{0,06} & \textbf{0,391} & \textbf{não} \\
\% de saídas por COVID (2020-21) & 9,8 & 8,7 & 5,2 & 0,04 & 0,011 & \textbf{invertida} \\
\bottomrule
\end{tabular}
\end{table}

\textbf{A diferença entre a metade de cima e a metade de baixo da tabela é o
achado desta subseção.} Em cima estão os indicadores de tamanho e de
operação --- leitos, saídas, diárias, faturamento, ocupação --- e neles o
degrau da escada explica de 23\% a 43\% do que separa um hospital do outro.
Embaixo estão os indicadores de resultado, e ali o mesmo degrau explica
pouco: na mortalidade, apenas 6\%.

Traduzido para quem decide: \textbf{saber em que degrau um hospital está
permite prever com folga o tamanho dele e o ritmo com que ele opera, e quase
não permite prever quantos dos seus pacientes morrem.} A escada organiza
capacidade instalada. Ela não organiza desfecho clínico.

\textbf{E na mortalidade a escada nem sequer sobe.} Os valores medianos são
4,3\% no primeiro degrau, 5,9\% no segundo e 5,8\% no terceiro --- o
terceiro fica \emph{abaixo} do segundo, e o teste diz que os dois são
indistinguíveis ($p_{23}$ de 0,391, muito acima de 0,05). O que separa a
mortalidade, portanto, é uma coisa só: \textbf{ter ou não ter terapia
intensiva}. Entre os hospitais que têm, subir na escala de complexidade não
mexe no resultado. É a mesma conclusão que o quadro anterior já sugeria pelo
empate dos números, agora com o teste feito e a diferença medida.

\subsubsection*{Três padrões que só aparecem com o painel largo}

\textbf{Primeiro: a produtividade por leito satura em E2.} As saídas por
leito SUS vão de 43,8 em E1 para 61,9 em E2 e \emph{recuam} para 59,0 em E3,
com $p$ de 0,736 entre os dois últimos. O ordenamento se quebra aqui, e
quebra num lugar interessante: toda a estrutura adicional de E3, que é mais
que o dobro de leitos, \textbf{não compra mais giro por leito}. O que E3
entrega a mais é volume absoluto e complexidade do caso, não eficiência
física do leito.

\textbf{Segundo: a composição público-privada não segue a escada.} A
proporção de leitos SUS sobre o total é de 80,5\% em E1, cai para 63,7\% em
E2 e volta a 81,2\% em E3. O estrato do meio é o mais misto, e é o único em
que um terço dos leitos atende fora do SUS. A não-monotonicidade é modesta
em tamanho ($\eta^2$ de 0,08) mas é sistemática, e vale como advertência: o
degrau intermediário não é ``E3 menor'', é um objeto de natureza diferente.

\textbf{Terceiro: a pandemia entrou pela ponta de baixo da escada.} A
fração de saídas por COVID no biênio 2020--2021 é de 9,8\% em E1, 8,7\% em
E2 e 5,2\% em E3, e é o \textbf{único ordenamento invertido} de todo o
levantamento. Os hospitais de menor complexidade absorveram
proporcionalmente mais COVID que os de alta, o que é coerente com uma rede
que reservou a capacidade de alta complexidade para o caso não pandêmico,
mas o dado aqui é descritivo e não sustenta afirmação sobre política de
regulação de leitos.

As três figuras a seguir são três leituras do mesmo corte: escala contra
intensidade tecnológica, a ocupação de internação separada nas suas duas
perguntas, e giro contra permanência. Cada uma tem uma marca por hospital, e
não por hospital-ano: o ponto é a mediana dos onze anos daquele hospital.

\begin{figure}[H]
\centering
\includegraphics[width=\textwidth]{figSEG_12_escala.png}
\caption{Escala e intensidade tecnológica. Cada ponto é um hospital: no
eixo horizontal, o total de leitos, em escala logarítmica; no vertical,
quanto desses leitos é de UTI. Os hospitais de E1, na base do gráfico, são
os que não têm UTI SUS.}
\label{fig:estruturauso}
\end{figure}

\comofeito{dispersão com uma marca por hospital (mediana dos onze anos) e
cor por estrato: leitos totais no eixo horizontal, em escala logarítmica, e
a proporção de leitos de UTI no vertical.}

\begin{figure}[H]
\centering
\includegraphics[width=\textwidth]{figSEG_12_duas_perguntas.png}
\caption{A ocupação de internação são duas perguntas: quão cheio o hospital
opera, e com que frequência ele passa da própria capacidade. A barra cinza
mistura as duas; a verde responde a primeira e a azul, a segunda.}
\label{fig:duasperguntas}
\end{figure}

\comoler{compare as barras cinza e verde de cada estrato: quando elas se
afastam, é porque a mediana bruta está sendo puxada pela frequência de
superlotação, não pelo nível de ocupação. Em E1 as duas coincidem; em E3 há
sete pontos entre elas.}

\comofeito{barras: para cada estrato, a mediana bruta da ocupação de
internação, a mediana calculada só sobre os hospitais-ano com ocupação até
100\%, e a proporção de hospitais-ano acima de 100\%.}

\textbf{Esta figura fecha o argumento da subseção seguinte com número
próprio.} A mediana bruta de ocupação vai de 41\% em E1 a 90\% em E3, mas
parte disso não é ocupação: é superlotação entrando no numerador. Separadas
as duas, o nível de ocupação abaixo de 100\% vai de 41\% a 83\% e a
frequência de superlotação vai de \textbf{1\% dos hospitais-ano em E1 para
9\% em E2 e 28\% em E3}. A distância entre a barra cinza e a verde é de zero
ponto em E1 e de sete pontos em E3. A separação formal das duas grandezas,
com as definições e a identidade que as liga, está na
Subseção~\ref{sub:lotacao}; toda linha de ``ocupação de internação'' desta
seção e das anteriores é a \textbf{lotação} ali definida, que pode passar
de 100\%, e não a taxa de ocupação, que não pode.

\begin{figure}[H]
\centering
\includegraphics[width=\textwidth]{figSEG_12_giro.png}
\caption{Giro contra permanência. Quanto mais à direita, mais saídas o
hospital produz por leito; quanto mais abaixo, menos tempo cada paciente
fica internado. A nuvem desce da esquerda para a direita porque as duas
coisas são a mesma: quem segura menos o leito, gira mais.}
\label{fig:giro}
\end{figure}

\comofeito{dispersão com uma marca por hospital (mediana dos onze anos) e
cor por estrato: saídas por leito SUS no eixo horizontal e tempo médio de
permanência no vertical.}

% A figura figSEG_13_estratos_por_ano.png saiu daqui em 04/08/2026, apanhada
% por João numa leitura de ponta a ponta. Ela era um painel de seis métricas
% ano a ano por estrato, e CINCO das seis já estavam publicadas, uma por
% página, nas Figuras das séries por estrato acima (figSEG_02_serie_*):
% mortalidade, TMP, faturamento por saída, ocupação de internação e fração
% de alta complexidade. Mesma mediana, mesma unidade, mesma faixa cinza de
% pandemia, mesma seção. A sexta métrica, saídas por leito SUS, já tinha o
% seu achado narrado no primeiro dos "Três padrões" acima, com o teste
% junto. A figura continua sendo gerada por descritivas_estratos_ampliada.py;
% o que saiu foi a republicação.
A leitura ano a ano das séries acima acrescenta duas coisas. O \textbf{TMP e
o faturamento por saída} correm em três faixas separadas e paralelas os onze
anos inteiros: são as métricas em que os estratos são mais estáveis como
grupos. Já a \textbf{mortalidade e as saídas por leito} mostram E2 e E3
trocando de posição repetidamente: na mortalidade as duas curvas se cruzam
\textbf{seis vezes} em onze anos, e nas saídas por leito o E2 fica acima do
E3 em \textbf{seis dos onze anos, inclusive nos três últimos}. O gradiente
estrutural é firme na estrutura e no preço unitário, e frouxo justamente
onde se mediria desempenho.

% O parágrafo "Onde ficam as tabelas completas" saiu daqui em 04/08/2026, a
% pedido de João: era um índice de nomes de arquivo de trabalho
% (tabSEG_20 a tabSEG_25) endereçado a quem tem o repositório aberto, e não
% ao leitor do documento. As seis tabelas continuam existindo, geradas por
% descritivas_estratos_ampliada.py, todas com a marca de tarja onde ela é
% devida. O que a subseção declarava e que vale manter — que esta parte
% mostra as duas pontas de um levantamento de 27 variáveis — já está dito na
% abertura da própria subseção.
\subsection{Os resultados em E3}

Esta é a comparação que interessa a quem decide, porque é a única feita
entre hospitais de fato parecidos: \textbf{125 hospitais de capacidade
instalada equivalente, 28 geridos por OSS e 20 pela administração Direta}.
Tudo o que segue compara essas duas redes dentro desse grupo, descontando
complexidade, porte e ano.

\begin{table}[H]
\centering
\caption{O que muda sob gestão por OSS entre hospitais de capacidade
equivalente (estrato E3), contra a administração Direta. A coluna
\textbf{Diferença} é o resultado; a \textbf{margem} diz entre que valores
ele pode estar, e quando ela cobre o zero a diferença não está
demonstrada.}
\label{tab:e3resultados}
\small
\begin{tabular}{llrr}
\toprule
Indicador & Diferença nas OSS & Margem (95\%) & Demonstrada? \\
\midrule
Tempo de permanência      & $\mathbf{-11{,}8\%}$ de tempo internado & $[-17{,}6\%;\ -5{,}6\%]$ & sim \\
Mortalidade geral         & menor, $-0{,}281$ na escala do modelo & $[-0{,}442;\ -0{,}121]$ & sim, com ressalva \\
Ocupação até 100\%        & $+27{,}0\%$ de leito ocupado & $[-1{,}8\%;\ +64{,}3\%]$ & não, por pouco \\
Ocupação de UTI           & $+25{,}5\%$ de uso da UTI & $[-18{,}4\%;\ +93{,}0\%]$ & não \\
Faturamento real/saída    & $-9{,}2\%$ por paciente & $[-25{,}4\%;\ +10{,}4\%]$ & não \\
Superlotação              & mais frequente, $+0{,}74$ na escala do modelo & $[-1{,}17;\ +2{,}64]$ & não \\
\bottomrule
\end{tabular}
\end{table}

\comofeito{cada linha é um modelo estimado no estrato E3, com efeitos
aleatórios correlacionados de Mundlak e erros-padrão agrupados por
hospital. \textbf{Diferença} é o efeito de OSS contra Direta traduzido para
a unidade de leitura; nas duas linhas em que aparece ``escala do modelo'' a
tradução para percentual não é possível sem induzir ao erro, porque o
modelo trabalha em logito. \textbf{Margem} é o intervalo de confiança de
95\%. \textbf{Demonstrada} responde se esse intervalo exclui o zero.}

\textbf{Dois resultados se sustentam entre hospitais comparáveis: o tempo
de permanência e a mortalidade.} O paciente fica \textbf{11,8\% menos tempo
internado} sob gestão por OSS, e a mortalidade sai menor. Os outros quatro
têm margem que cobre o zero e não autorizam afirmação em direção nenhuma,
inclusive a ocupação até 100\%, que passa perto.

\textbf{O faturamento não decide, e a tentação de fazê-lo decidir tem
nome.} A margem vai de 25,4\% de queda a 10,4\% de alta, ou seja, cobre o
zero: não se pode dizer que as OSS faturam menos por paciente em E3.

\textbf{Sobre a mortalidade, a ressalva é forte e precisa acompanhar o
número.} O resultado da tabela vale \textbf{com a complexidade do hospital
descontada}, e só assim. \textbf{A diferença bruta de mortalidade entre OSS
e Direta dentro de E3 não é citável}, e por isso não aparece neste
documento: parte dela é composição de casuística dentro do próprio estrato,
porque a Direta concentra hospitais na faixa de complexidade máxima, em que
as OSS não têm nenhum, e essa assimetria vem de dois hospitais.

\subsubsection*{Sobre o que estas estimativas repousam}

\textbf{Esta é a qualificação mais importante da seção, e ela vale para
todos os seis números da tabela.}

O que a tabela mede é a \emph{mudança}, e quem produz mudança são os
hospitais que trocaram de rede no período; os demais servem para apertar a
margem, não para determinar o resultado. Em E3 são \textbf{quatro}
hospitais, e eles pesam de forma muito desigual: dois converteram em 2025 e
trazem \textbf{um único ano} de operação depois da mudança. Na prática,
\textbf{tudo repousa sobre o Conjunto Hospitalar de Sorocaba e o Pérola
Byington}, os únicos com histórico longo depois da conversão.

E os dois trazem complicação conhecida. Em Sorocaba, a conversão foi
decidida em resposta a uma crise assistencial, ou seja, o hospital estava
piorando quando trocou de gestão, e a comparação antes e depois mistura o
efeito da nova gestão com a simples volta ao patamar normal. No Pérola
Byington, a conversão coincidiu com mudança de prédio. Soma-se a isso que
\textbf{os dois foram para a mesma organização social}, o Seconci-SP
(Subseção~\ref{sub:oqueforam}): o que se está medindo é o desempenho
daquele operador, e não do modelo OSS em geral. Dos outros três
conversores, dois foram para uma segunda organização, e de um deles não se
identificou nem a data exata da assunção nem a organização contratada.

A margem de erro não sabe nada disso. Ela mede a incerteza que vem do
tamanho da amostra, não a que vem de o experimento não ter sido desenhado.
\textbf{Não é um desenho limpo; é o desenho disponível}, e os números desta
seção devem circular com essa frase colada.

\subsubsection*{E1 e E2 não são interpretados}

Os coeficientes de E1 e E2 foram estimados e ficaram \textbf{tarjados}.
Não são narrados aqui, e o motivo é estrutural, não de cautela retórica:
\textbf{E1 não tem nenhum hospital que trocou de rede} e E2 tem um só. Sem
conversores, o coeficiente de OSS deixa de ser efeito e passa a ser
\textbf{comparação entre grupos distintos de hospitais}, a mesma coisa que
a segmentação existe para não fazer.

Um exemplo do que isso produz: o TMP em E1 aparece 72,9\% menor nas OSS. Esse
número é o contraste entre \textbf{dois hospitais OSS e um da Direta}. Não é
achado; é ruído com tarja, e não recebe prosa.

O \textbf{Privado não é interpretado em nenhum estrato}, aqui como no resto
do documento: são 3 hospitais no painel inteiro, e nenhum em E3.

% A subseção "A eficiência técnica, medida dentro do estrato" saiu daqui em
% 04/08/2026, apanhada por João: ela usava análise envoltória de dados,
% fronteira, correção de viés e regressão de segunda etapa DUAS SEÇÕES ANTES
% de o documento explicar o que é qualquer uma dessas coisas. O leitor
% encontrava o resultado de um método que ainda não tinha sido apresentado.
% Foi movida inteira para a Seção 16 (`sec15_fase2_eficiencia.tex`), logo
% depois do DEA global, onde o leitor já sabe o que é fronteira e o que
% significa estar perto dela. A figura abaixo, que é de coeficiente por
% estrato e de superlotação por categoria, NÃO era da subseção movida e
% permanece aqui.
\begin{figure}[H]
\centering
\subfig{figSEG_04_coef_oss_por_estrato.png}{Efeito por estrato}
\subfig{figSEG_05_sobrecarga_por_categoria.png}{Superlotação por categoria}
\caption{No painel de cima, o coeficiente de OSS em cada estrato; as barras em
cinza são as tarjadas, sem identificação within. No de baixo, a frequência de
superlotação por categoria de gestão dentro de cada estrato.}
\label{fig:e3coef}
\end{figure}

\comofeito{no painel de cima, um gráfico de pontos com barras de intervalo: o
eixo vertical traz o estrato e o horizontal, o \textbf{efeito estimado da
gestão por OSS contra a administração Direta} naquele estrato, com o
intervalo de 95\% de erros-padrão agrupados por hospital. As barras em
cinza são E1 e E2, tarjadas porque não têm conversores suficientes para
identificar efeito \emph{within}: o ponto ali é contraste entre grupos
distintos de hospitais, não efeito de mudança de gestão. No de baixo, um
gráfico de barras: para cada estrato, a \textbf{proporção de
hospitais-ano em que a ocupação de internação passa de 100\%}, uma barra
por categoria de gestão. A unidade é o hospital-ano.}

\subsection{A ocupação de internação eram duas perguntas, não uma}
\label{sub:decomposicao}

% A parte DEFINICIONAL desta subseção — por que a ocupação passa de 100% e a
% distinção entre lotação e taxa de ocupação — foi movida para a Seção 2 em
% 04/08/2026, por decisão de João. O motivo é de ordem de leitura: ela governa
% toda taxa de ocupação do documento, inclusive as das Seções 4, 5 e 11, e
% ficava setenta páginas depois do primeiro uso. O que permanece aqui é o que
% de fato depende do estrato E3, que só existe nesta seção.
O que segue é o resultado que mais muda a leitura da ocupação de internação,
e ele parte da distinção já estabelecida na
Subseção~\ref{sub:lotacao}: o que o documento publica como ``ocupação'' é a
\textbf{lotação}, que pode passar de 100\%, e ela se separa em \textbf{leito
de fato ocupado} mais \textbf{transbordo}. Aqui essa separação é aplicada à
comparação entre as redes, dentro do único grupo em que ela é viável.

\subsubsection*{A separação, escrita com precisão}

Sendo $D_{it}$ as diárias de internação do hospital $i$ no ano $t$ e
$L_{it}$ os seus leitos SUS cadastrados, as três grandezas se definem assim:

\begin{equation*}
  \underbrace{\ell_{it} = \frac{D_{it}}{365\,L_{it}}}_{\text{lotação}}
  \qquad
  \underbrace{\tau_{it} = \min(\ell_{it},\,1)}_{\text{taxa de ocupação}}
  \qquad
  \underbrace{e_{it} = \max(\ell_{it}-1,\,0)}_{\text{excesso}}
\end{equation*}

\noindent acompanhadas do indicador de superlotação
$s_{it} = \mathbf{1}\{\ell_{it} > 1\}$, que vale um quando o hospital
demanda mais leito-dia do que tem habilitado. As três se ligam por uma
\textbf{identidade exata}, verificada no painel sem resíduo:

\begin{equation}
  \ell_{it} \;=\; \tau_{it} \;+\; e_{it}.
  \label{eq:lotacao}
\end{equation}

\textbf{A variável que este documento imprime como ``ocupação de
internação'' é $100\,\ell_{it}$, e não $100\,\tau_{it}$}, conforme já
declarado na Subseção~\ref{sub:lotacao}. Quando $\ell \leq 1$, lotação e
taxa de ocupação coincidem e o excesso é zero; é só acima de 100\% que elas
se separam, e é exatamente ali que mora o achado desta subseção.

Isto também explica, em uma linha, por que a Seção~\ref{sec:fase2} modela a
ocupação de internação \textbf{em duas partes} e não como contínua aparada:
aparar em 100\% equivale a jogar fora $e_{it}$, isto é, a estimar $\tau$ e
chamá-lo de $\ell$. O modelo em duas partes estima $s$ e $\tau$
separadamente, que são os dois termos da identidade.

\comoler{à esquerda, a distribuição da lotação no painel inteiro: a parte
verde é onde a lotação ainda é uma taxa de ocupação, e a parte alaranjada,
à direita da linha tracejada, é a superlotação. À direita, a
Equação~\ref{eq:lotacao} desenhada: a altura total de cada barra é a
lotação média da categoria, a parte de baixo é a taxa de ocupação e a
faixa alaranjada no alto é o excesso. Compare a altura da faixa alaranjada
com a frequência escrita dentro da barra: são coisas diferentes.}

\begin{figure}[H]
\centering
\includegraphics[width=\textwidth]{figSEG_14_decomposicao_lotacao.png}
\caption{A lotação do painel separada nas suas duas parcelas. À esquerda,
onde a lotação deixa de ser taxa de ocupação; à direita, a identidade
$\bar\ell = \bar\tau + \bar e$ dentro do estrato E3, o único em que a
comparação entre redes é viável.}
\label{fig:decomposicaolotacao}
\end{figure}

\comofeito{à esquerda, histograma da lotação $\ell$ (em \%) sobre os 3.179
hospitais-ano, em classes de 5 pontos, aparado em 250 apenas para a vista
(a cauda vai além, e nenhuma observação é descartada do cálculo). À
direita, barras empilhadas: para cada categoria de gestão dentro de E3, a
\textbf{média} de $\tau$ na base e a \textbf{média} de $e$ no topo, cuja
soma é a média de $\ell$ impressa acima da barra. A unidade é o
hospital-ano. O Privado fica fora do painel direito pelo tamanho (3
hospitais no painel inteiro, nenhum em E3).}

\subsubsection*{As duas perguntas, separadas}

Se o excesso mede superlotação, então a média de ocupação mistura duas
coisas distintas: \emph{com que frequência} o hospital entra em
superlotação, e \emph{quão cheio} ele opera quando não está superlotado. O
indicador foi separado nessas duas partes.

\begin{table}[H]
\centering
\caption{As duas perguntas da ocupação de internação em E3, por categoria.
A última coluna é a mediana bruta da lotação, que mistura as duas.}
\label{tab:duaspartes}
\small
\begin{tabular}{lrrr}
\toprule
Categoria em E3 & \% de anos em & Ocupação mediana & Mediana bruta \\
                & superlotação  & quando $\leq$100 & (conflaciona) \\
\midrule
Direta            &  2,9 & 67,9 & 68,5 \\
Filantrópico      & 28,7 & 83,5 & 89,9 \\
OSS               & 35,5 & 89,0 & 95,9 \\
Público Municipal & 36,4 & 86,5 & 94,5 \\
\bottomrule
\end{tabular}
\end{table}

A separação mostra que \textbf{as duas coisas andam juntas, mas não são a
mesma}. A distância entre a mediana bruta e a mediana condicional é de 0,6
ponto na Direta e de 6,9 pontos nas OSS: \textbf{parte do que aparecia como
``OSS opera mais cheio'' era, na verdade, ``OSS entra em superlotação com
muito mais frequência''}.

\subsubsection*{Quanto da vantagem das OSS é leito cheio, e quanto é transbordo}

A pergunta agora é de tamanho: das duas coisas que estavam somadas, qual
pesa mais na distância entre as redes? A conta se faz com a média em vez da
mediana, por um motivo prático --- somar medianas não dá a mediana da soma,
e é justamente a soma que precisa ser repartida aqui.

Entre os hospitais comparáveis, as OSS operam a \textbf{93,0\%} da
capacidade cadastrada contra \textbf{69,4\%} da administração Direta. São
\textbf{23,5 pontos} de distância, e eles se repartem assim:

\begin{equation*}
  \underbrace{23{,}5}_{\text{distância total}}
  \;=\;
  \underbrace{21{,}5}_{\text{leito de fato ocupado }(91\%)}
  \;+\;
  \underbrace{2{,}0}_{\text{transbordo }(9\%)}
\end{equation*}

\textbf{Nove por cento da distância é transbordo; noventa e um por cento é
leito de fato ocupado.} O número corrige duas leituras erradas ao mesmo
tempo, e as duas circulam.

A primeira é a leitura ingênua, de que o indicador bruto media só o quanto o
hospital opera cheio. Não media: quando a superlotação sai da conta, a
distância entre as duas redes encolhe --- de 27,4 para 21,1 pontos na
mediana ---, ou seja, \textbf{quase um quarto do que se via era frequência
de transbordo, e não nível de ocupação}. A segunda leitura errada é a
oposta, a de que a vantagem das OSS \emph{seria} superlotação disfarçada.
Também não é: tirando todo o transbordo, as OSS continuam operando bem mais
cheias que a Direta, por uma margem que a Direta não chega perto de
alcançar. \textbf{A diferença de ritmo entre as duas redes sobrevive
inteira; o que a separação corrige é o tamanho dela, não a sua existência.}

A mesma conta arruma as outras duas redes de um jeito que o número bruto
escondia. O Público Municipal parece empatar com as OSS (92,6\% contra
93,0\%), mas chega lá por outro caminho: \emph{mais} transbordo (3,2 contra
2,1 pontos) e \emph{menos} leito ocupado (89,4 contra 90,9). O Filantrópico
é o que mais transborda dos quatro (4,0 pontos) e, ainda assim, o segundo
com menos leito ocupado (84,7). A regra que emerge é útil e vale anotar:
\textbf{quanto maior a parcela de transbordo, menos ela fala de hospital
operando cheio, e mais de hospital com leito de menos para a demanda que
bate na porta}.

\subsubsection*{Superlotar mais vezes não é superlotar mais fundo}

Faltou separar duas perguntas que também vinham somadas: \emph{com que
frequência} o hospital passa da capacidade, e \emph{quanto} ele passa quando
passa. São coisas diferentes, e a rede que mais faz uma não é a que mais faz
a outra.

Primeiro, o tamanho típico do transbordo, que é menor do que a palavra
``superlotação'' sugere: nos anos em que um hospital passa de 100\%, ele
passa em geral para \textbf{106,9\%}, e em nove de cada dez casos fica
abaixo de 120\%. Não é hospital operando ao dobro da capacidade; é hospital
operando com alguns leitos a mais do que os cadastrados.

Segundo, e aqui a ordem das redes \textbf{se inverte}. As OSS são a rede que
\emph{mais vezes} transborda (35,5\% dos anos) e, ao mesmo tempo, a que
transborda \emph{menos fundo}: 104,6\% nos anos em que passa, contra 107,1\%
no Público Municipal e 108,9\% no Filantrópico, cujo pior caso chega a
236,6\%. A administração Direta aparece com 100,5\%, mas sobre apenas
\textbf{seis} hospitais-ano, e por isso o número fica tarjado e não entra na
comparação.

\textbf{Isso não contradiz o achado anterior; ele o completa.} A pressão
sobre as OSS é \emph{frequente e rasa} --- o hospital passa da capacidade
quase todo ano, e passa pouco. A do Filantrópico é \emph{rara e profunda}
--- passa poucas vezes, e quando passa, passa muito. São dois problemas de
gestão diferentes, que pedem respostas diferentes, e o indicador de ocupação
sozinho não distinguia um do outro porque somava os dois num número só.

\comoler{à esquerda, olhe a \emph{altura das caixas}: a das OSS é a mais
baixa das três não tarjadas, apesar de as OSS serem a rede que mais vezes
passa da capacidade. À direita, cada ponto é uma rede dentro de um degrau:
quanto mais à direita, mais vezes ela transborda; quanto mais alto, mais
fundo. Se as duas coisas fossem a mesma, os pontos formariam uma diagonal.
Eles não formam, e é esse o recado da figura.}

\begin{figure}[H]
\centering
\includegraphics[width=\textwidth]{figSEG_16_profundidade_excesso.png}
\caption{Frequência e profundidade da superlotação são dimensões
independentes. As OSS superlotam com o dobro da frequência do
Filantrópico e com quatro pontos a menos de profundidade mediana.}
\label{fig:profundidadeexcesso}
\end{figure}

\comofeito{à esquerda, boxplots da lotação \textbf{restritos aos
hospitais-ano em que ela passa de 100}, um por categoria de gestão dentro de
E3; a linha branca é a mediana, a caixa vai do primeiro ao terceiro quartil
e os pontos soltos são observações extremas. O $n$ de cada caixa está no
eixo, e a Direta sai tarjada por ter menos de dez observações. À direita,
dispersão com uma marca por categoria em cada estrato: no eixo horizontal a
\textbf{frequência} (\% de hospitais-ano acima de 100) e no vertical a
\textbf{profundidade} (lotação mediana condicional a passar de 100). Marca
cheia e grande é E3, marca clara e pequena é E2; E1 fica fora por ter quase
nenhuma superlotação.}

\subsubsection*{Quantos hospitais superlotam, e por quantos anos}

Há ainda uma terceira dimensão, que nem a frequência agregada nem a
profundidade capturam: a superlotação se concentra em poucos hospitais
crônicos ou se espalha por muitos hospitais em anos isolados? A resposta
importa porque as duas leituras têm implicações opostas para a natureza do
excesso, e a pergunta se responde contando hospitais, não observações.

\textbf{No painel inteiro, o padrão é episódico.} Dos 289 hospitais, 123
passam de 100\% em ao menos um ano, e \textbf{166 nunca passam}. Entre os
123 afetados, a mediana é de apenas \textbf{três anos} dos onze, 57 têm um
ou dois anos e \textbf{só três hospitais superlotam nos onze anos}. Uma
capacidade cronicamente sub-registrada produziria excesso permanente, e não
é isso que se vê: o retrato geral é o de picos de demanda.

\textbf{Mas o padrão geral não vale para todas as redes, e a diferença é
grande.} Nas OSS, 20 dos 33 hospitais (60,6\%) superlotam alguma vez, e
entre esses a mediana é de \textbf{5,5 anos dos onze}: metade da década.
No Público Municipal o alcance é ainda maior (39 de 56, ou 69,6\%) com
recorrência menor (mediana de 3 anos). No Filantrópico e na Direta a
superlotação atinge cerca de um terço dos hospitais (32,2\% e 30,8\%) e
dura, na mediana, dois anos e um ano. \textbf{Nas OSS a superlotação não é
episódica: é o modo normal de operar de boa parte da rede.}

Vale registrar o que isso faz com um argumento anterior. O padrão episódico
era uma das evidências de que o excesso mede sobrecarga real e não
sub-registro de leito, e ele continua valendo \textbf{para o painel como um
todo}. Para as OSS, especificamente, ele deixa de valer: um hospital que
opera acima de 100\% em metade dos anos é compatível tanto com pressão
assistencial persistente quanto com uma capacidade habilitada que ficou
defasada da capacidade operada. \textbf{O mecanismo que sustenta a leitura de
sobrecarga é o das definições do CNES e do SIH}, estabelecido no começo desta
subseção, e não o padrão temporal; este apenas o acompanhava, e para uma rede
deixa de acompanhar.

\comoler{à esquerda, a barra cinza da posição zero são os hospitais que
nunca superlotam, e é de longe a maior: o grosso da rede não superlota
nunca. Entre os que superlotam, as barras caem depressa, o que é o padrão
episódico. À direita, quanto mais à direita a rede está, mais dos seus
hospitais superlotam alguma vez; quanto mais alto, por mais anos. As OSS
estão altas.}

\begin{figure}[H]
\centering
\includegraphics[width=\textwidth]{figSEG_17_persistencia_superlotacao.png}
\caption{Alcance e recorrência da superlotação. Episódica no painel como um
todo (mediana de três anos entre os afetados), ela é recorrente nas OSS,
onde atinge 60,6\% dos hospitais por uma mediana de 5,5 dos onze anos.}
\label{fig:persistencia}
\end{figure}

\comofeito{à esquerda, contagem de hospitais (CNES) por número de anos, de
zero a onze, em que a lotação passa de 100\%; a unidade é o hospital, não o
hospital-ano, e a barra do zero está em cinza por não ser superlotação. À
direita, uma marca por categoria de gestão: no eixo horizontal o
\textbf{alcance} (\% dos hospitais da rede que superlotam ao menos uma vez)
e no vertical a \textbf{recorrência} (mediana de anos entre os que
superlotam). Entre parênteses, quantos hospitais da rede são afetados sobre
o total dela. O Privado fica fora pelo tamanho.}

\subsubsection*{As três leituras ano a ano}

Decompostas as três grandezas, a série temporal deixa de ser ambígua: dá
para ver separadamente o que aconteceu com a lotação, com a ocupação e com
a frequência de superlotação em cada rede.

Três coisas aparecem. A primeira é que \textbf{a distância entre a Direta e
as demais redes é estrutural e atravessa a década inteira}, nas três
leituras, e não é artefato de um ano ou do biênio pandêmico. A segunda é que
\textbf{a pandemia esvaziou a Direta e não esvaziou as OSS}: em 2020 a taxa
de ocupação média da Direta em E3 cai a 60,6\%, o menor valor de toda a
série, enquanto a das OSS fica em 88,2\%, praticamente no seu patamar
habitual. A terceira é que \textbf{a superlotação cresce no fim da série},
e não nela: nas OSS a frequência vai de 33,3\% em 2015 a 50,0\% em 2024, e
no Público Municipal de 35,5\% a 51,6\%, enquanto a Direta permanece em zero
em sete dos onze anos. O aperto de leito da rede paulista, nestes dados, é
mais um fenômeno de 2024 e 2025 do que de 2020 e 2021.

\comoler{compare os três painéis lendo a linha da Direta, que é a mais
baixa em todos: no painel (b) ela cai fundo em 2020, e no (c) ela encosta no
zero e fica. As três leituras contam partes diferentes da mesma história, e
só a do meio é uma taxa de ocupação no sentido próprio da palavra.}

\begin{figure}[H]
\centering
\includegraphics[width=\textwidth]{figSEG_15_lotacao_por_ano.png}
\caption{Lotação, taxa de ocupação e frequência de superlotação ano a ano
no estrato E3, por categoria de gestão. A faixa cinza é o biênio
pandêmico.}
\label{fig:lotacaoano}
\end{figure}

\comofeito{três painéis com o mesmo eixo horizontal, o ano de 2015 a 2025,
e uma linha por categoria de gestão dentro de E3. No painel (a), a média
anual da lotação $\ell$; no (b), a média anual da taxa de ocupação $\tau$,
que é a mesma grandeza com teto em 100; no (c), a \% de hospitais-ano da
categoria em superlotação, isto é, a média anual de $s$. A unidade é o
hospital-ano e a faixa cinza marca 2020 e 2021. O Privado fica fora por não
ter hospital em E3.}

\subsubsection*{O que os dados não decidem}

Aqui é preciso ser explícito sobre um limite, porque a tabela acima convida a
uma conclusão que ela não sustenta.

\textbf{A diferença de superlotação entre as redes é grande e bem medida}:
2,9\% dos anos na Direta contra 35,5\% nas OSS, dentro do mesmo estrato
estrutural. \textbf{Isso é descrição, e é sólida.}

O que \emph{não} se pode dizer é que a gestão por OSS \emph{cause} essa
diferença. A estimativa que tentaria medir isso (a parte A da tabela de
resultados, $+0{,}74$ no logito) repousa sobre os hospitais que trocaram de
rede, e ali o material é escasso ao ponto de não decidir nada. \textbf{São
quatro hospitais no painel inteiro que sustentam essa estimativa}, e o exame
detalhado é ainda menos animador: \textbf{apenas dois deles registram
superlotação do lado posterior à conversão}, o lado que efetivamente
identifica o efeito. Nos outros dois, toda a superlotação observada é
anterior à troca de gestão.

O intervalo resultante vai de $-1{,}17$ a $+2{,}64$ no logito: ele cobre
desde uma redução enorme até um aumento enorme. \textbf{Uma estimativa nula
sobre quatro hospitais não é evidência de que não há efeito: é evidência de
que estes dados não têm potência para detectá-lo.} Não se deve escrever que
a gestão não afeta a superlotação, nem que afeta.

\subsection{O que a segmentação acrescenta}

Três coisas, em ordem de importância.

\textbf{Primeira: ela delimita onde há evidência.} Dos três estratos,
\textbf{apenas um} sustenta comparação entre redes. Isso não é limitação
introduzida pela segmentação: é limitação que já existia e que a análise
agregada escondia, ao somar num único contraste hospitais que não têm
correspondente na rede oposta.

\textbf{Segunda: ela confirma que a partição é real.} O gradiente monotônico
em cinco dos seis indicadores mostra que os três degraus capturam diferença
estrutural de fato, e não um recorte arbitrário.

\textbf{Terceira: ela separou duas perguntas que estavam somadas} no
indicador de ocupação, e mostrou que a distância entre as redes se distribui
diferentemente entre elas. Com as definições da Subseção~\ref{sub:lotacao},
a separação deixa de ser qualitativa: dos 23,5 pontos que separam a lotação
média das OSS da administração Direta em E3, \textbf{21,5 são taxa de
ocupação e 2,0 são excesso}. A conflação existia e era real, mas explicava
9\% do contraste, não o contraste.

\textbf{E uma quarta, que o levantamento ampliado da
Subseção~\ref{sub:ampliada} acrescenta e que é de natureza diferente das
três:} a escada de complexidade organiza fortemente a \emph{estrutura} dos
hospitais e quase não organiza o \emph{desfecho}. O estrato explica de 23\%
a 43\% da variância entre hospitais em leitos, volume, faturamento e
ocupação, e 6\% na mortalidade, onde nem sequer ordena. Isso é uma
advertência sobre o próprio instrumento: segmentar por capacidade instalada
resolve o problema de comparar coisas incomparáveis, mas não deve ser lido
como se o degrau explicasse a qualidade do que o hospital entrega.

O que a segmentação \emph{não} faz é resolver o problema de desenho. Em E3,
como no painel inteiro, a identificação causal continua repousando sobre
poucos hospitais convertidos, com confundidores documentados e um único tipo
de operador.
